\section*{Chapter \ref{ch:rs:measuring-market-making-and-execution} --- Measuring Market Making and Execution}

\begin{solution}{exo:rs:measuring-market-making-and-execution:1}
\omterm{def:rs:measuring-market-making-and-execution:pnl}{Spread capture} $100.00 - 99.99 = +1$ tick; adverse selection $99.97 - 100.00 = -3$ ticks; inventory $100.02 - 99.97 = +5$ ticks; the total, $100.02 - 99.99 = +3$ ticks,
is their sum.
\end{solution}

\begin{solution}{exo:rs:measuring-market-making-and-execution:2}
$8\,144/18\,000 = 45\%$. Quotes are filled most readily when they are wrong (the price is about to move through them), so a higher \omterm{def:rs:measuring-market-making-and-execution:fill}{fill rate} can mean more adverse
selection; the filtered quoter fills 20\% and loses less.
\end{solution}

\begin{solution}{exo:rs:measuring-market-making-and-execution:3}
$0.43 - 1.10 = -0.67$ tick a share, the curve's value at twenty seconds.
\end{solution}

\begin{solution}{exo:rs:measuring-market-making-and-execution:4}
For each fill the three terms telescope to $m(T) - p$ whatever $H$ is. $H$ moves the boundary between adverse selection and inventory: too short, and the adverse
move still unfolding after $H$ is booked as inventory (adverse selection looks smaller, inventory worse); too long, and chosen inventory risk is booked as adverse
selection.
\end{solution}

\begin{solution}{exo:rs:measuring-market-making-and-execution:5}
Delay $80 \times 0.20 = 16$; execution $80 \times (50.40 - 50.20) = 16$; opportunity $20 \times (51.00 - 50.00) = 20$; fees $0.80$; shortfall \$52.80, which is the paper
portfolio's gain of \$100 minus the real one's \$47.20.
\end{solution}

\begin{solution}{exo:rs:measuring-market-making-and-execution:6}
The loss fell by \$4\,240. At the touch quoter's loss per lot (\$0.85), the filtered quoter's 3\,309 lots would have lost \$2\,814: fewer fills account for \$4\,111
(97\%), better fills for \$129.
\end{solution}

\begin{solution}{exo:rs:measuring-market-making-and-execution:7}
$6.5 \times 252 = 1\,638$ hours, 136.5 times the twelve simulated: \omterm{def:rs:measuring-market-making-and-execution:pnl}{spread capture} $+\$480\,000$, adverse selection $-\$1.23$ million, inventory $-\$145\,000$, fees
$-\$56\,000$, total $-\$945\,000$. The scaling assumes every hour of the year is like the simulated ones: the same activity, no open or close, no news, no change of
competitors.
\end{solution}

\begin{solution}{exo:rs:measuring-market-making-and-execution:8}
One order's shortfall against arrival is dominated by the market's move over its life: the chapter's seed-61 order beat arrival by 8.69 ticks because the price fell,
and the twelve orders' average of $-2.00$ has a standard error of 2.73. Measure the price paid against the mid at each fill, average over many orders with standard
errors, condition on the market's move, and compare algorithms by randomised experiment.
\end{solution}

\begin{solution}{pb:rs:measuring-market-making-and-execution:1}
\begin{enumerate}
\item 0.43 tick against the mid, 0.19 against the \omterm{def:rs:order-book-features:micro}{microprice}: the \omterm{def:rs:order-book-features:micro}{microprice} leans towards the next move, so part of the apparent half spread is already lost at
the fill.
\item Informed: $-1.85$ and $-2.71$. Uninformed: $+0.18$ and $+0.52$.
\item 42\%.
\item Twenty seconds: from there on the curve stays within two standard errors of its five-minute value.
\item Touch quoter 18\,000 posted, 8\,144 filled, 45\%; filtered 16\,391, 3\,309, 20\%.
\item It raises the spread captured (0.52 against the mid, 0.44 against the \omterm{def:rs:order-book-features:micro}{microprice}) and the twenty-second mark-out ($-0.53$ against $-0.67$); it does not reduce
the share of informed counterparties (40\% against 42\%).
\item $-\$0.85$ and $-\$0.81$.
\item \omterm{def:rs:measuring-market-making-and-execution:pnl}{Spread capture} \$3\,519, adverse selection $-\$8\,977$, inventory $-\$1\,060$, fees $-\$407$, total $-\$6\,925$.
\item \$1\,734, $-\$3\,483$, $-\$771$, $-\$165$, total $-\$2\,685$.
\item Per share they are the curve's value at the fill (0.43) and its fall to the horizon ($-1.10$); their sum is the twenty-second mark-out.
\item \texteuro1.55 a trade on the spread net of fees, a positioning loss of \texteuro0.68 (a gain on positions under five seconds, a loss on longer ones), \texteuro0.88
gross.
\item 28\,700 shares filled; delay 0.00, execution $-7.83$, opportunity $-0.91$, fees 0.05, shortfall $-8.69$ ticks a share; \omterm{def:rs:measuring-market-making-and-execution:costs}{VWAP slippage} $-1.10$.
\item Shortfall $-2.00$ (2.73); delay $+2.83$ (0.94); execution $-4.68$ (2.09), of which drift $-4.78$ (2.07) and price paid $+0.10$ (0.03); opportunity $-0.20$ (0.20);
fees 0.05.
\item The price paid against the mid at each fill, $+0.10$ tick a share; the rest is the market.
\item $+0.25$ tick (0.24): the mid moved $-9.71$ with the order and $-9.96$ without. The simulated efficient price ignores the order, so no impact lasts.
\item \textbf{Named result.} Scaled to a year: \omterm{def:rs:measuring-market-making-and-execution:pnl}{spread capture} $+\$480\,000$, adverse selection $-\$1.23$ million, inventory $-\$145\,000$, fees $-\$56\,000$, total
$-\$945\,000$; the mark-outs settle at twenty seconds.
\item Meet fewer informed traders (a fair-value signal that anticipates the efficient price, not the queue alone) or quote wider when they are active; inventory is
not the problem.
\item By randomised experiment on the firm's flow, on the price paid against the mid and on mark-outs, with standard errors (chapter~21).
\item The benchmark, the decomposition, the number of orders, the standard errors and what was conditioned on.
\item What happens to the price after each fill, averaged over fills and horizons: how much of the spread the fills keep.
\end{enumerate}
\end{solution}

\begin{solution}{iq:rs:measuring-market-making-and-execution:1}
The fill's P\&L per share against the reference price at a later horizon, $s(m(t+h) - p)$; use a curve of horizons from the fill to the point where it settles
(twenty seconds for the chapter's quoter), against both the mid and the \omterm{def:rs:order-book-features:micro}{microprice}.
\end{solution}

\begin{solution}{iq:rs:measuring-market-making-and-execution:2}
\omterm{def:rs:measuring-market-making-and-execution:pnl}{Spread capture} (the reference at the fill minus the price), adverse selection (the reference's move to a horizon $H$), inventory (after $H$), fees; exact by
construction. Choose $H$ where the \omterm{def:rs:measuring-market-making-and-execution:curve}{mark-out curve} settles, and report the sensitivity to it.
\end{solution}

\begin{solution}{iq:rs:measuring-market-making-and-execution:3}
It earns the spread and gives back less than half of it on positions, a profit of \texteuro0.88 a trade, many times a day; Menkveld found the positions profitable
under five seconds and costly beyond, so its edge is speed of unwinding and its main risk is being left holding inventory.
\end{solution}

\begin{solution}{iq:rs:measuring-market-making-and-execution:4}
Not necessarily: a \omterm{def:rs:measuring-market-making-and-execution:fill}{hit ratio} rises when the quotes become more generous than the competition's, which is when clients are most informed about the price. Look at
the mark-outs of the trades won.
\end{solution}

\begin{solution}{iq:rs:measuring-market-making-and-execution:5}
Randomise orders (or stock-days) between them, measure costs anchored at each fill (price paid against the mid, mark-outs) and the shortfall conditioned on the
market's move, and use CUPED-style adjustments and enough orders for the standard error (chapter~21).
\end{solution}

\begin{solution}{iq:rs:measuring-market-making-and-execution:6}
Inputs: every fill with side, price, quantity, time stamps (exchange and local), order and strategy ids, counterparty or flow type where known; the market data to
rebuild the mid and \omterm{def:rs:order-book-features:micro}{microprice}. Fixed horizons; groupings by strategy, venue, counterparty, time of day, size. Checks: decomposed P\&L equal to the accounting
P\&L to the cent; fill counts equal to the exchange's; references taken strictly before each fill.
\end{solution}
